% Иллюстрация к выводу распределения Больцмана:
% мысленный параллелепипед газа dx·dy·dz, баланс сил давления на гранях
% x и x+dx и потенциальной силы F (равновесие вдоль оси x).
% Компиляция: pdflatex volume-element-figure.tex
\documentclass[tikz,border=8pt]{standalone}
\usepackage[T2A]{fontenc}
\usepackage[utf8]{inputenc}
\usepackage[russian]{babel}
\usetikzlibrary{arrows.meta,calc}

\begin{document}
\begin{tikzpicture}[
    >={Stealth[length=2.4mm]},
    % --- стили в духе предыдущих рисунков ---
    axis/.style   = {->, semithick},
    helper/.style = {densely dashed, gray!70},
    presf/.style  = {->, line width=1.1pt, red!70!black},
    potf/.style   = {->, very thick, blue!55!black},
    comp/.style   = {->, semithick, blue!55!black!75},
  ]

  % ================= вершины бруска =================
  % передняя грань (ближе к наблюдателю)
  \coordinate (A)   at (3.2,1.5);     % перед-низ-лево
  \coordinate (Ax)  at (6.2,1.5);     % перед-низ-право
  \coordinate (Axy) at (6.2,3.7);     % перед-верх-право
  \coordinate (Ay)  at (3.2,3.7);     % перед-верх-лево
  % задняя грань: сдвиг "вглубь" на (1.265,0.806)
  \coordinate (B)   at (4.465,2.306);
  \coordinate (Bx)  at (7.465,2.306);
  \coordinate (Bxy) at (7.465,4.506);
  \coordinate (By)  at (4.465,4.506);

  % ================= оси =================
  \draw[axis] (0,0) -- (9.0,0) node[below right=-2pt] {$x$};
  \draw[axis] (0,0) -- (0,4.5) node[above] {$y$};
  \draw[axis] (0,0) -- (-1.85,-1.18) node[below left=-2pt] {$z$};

  % проекции граней x и x+dx на ось
  \draw[helper] (3.2,1.44) -- (3.2,0.07);
  \draw[helper] (6.2,1.44) -- (6.2,0.07);
  \draw (3.2,-0.07) -- (3.2,0.07);
  \draw (6.2,-0.07) -- (6.2,0.07);
  \node[below] at (3.2,-0.10) {$x$};
  \node[below] at (6.2,-0.10) {$x+\mathrm{d}x$};

  % ================= брусок =================
  % лёгкая заливка видимых граней
  \fill[gray!5]  (A) -- (Ax) -- (Axy) -- (Ay) -- cycle;   % передняя
  \fill[gray!12] (Ay) -- (Axy) -- (Bxy) -- (By) -- cycle; % верхняя
  \fill[gray!20] (Ax) -- (Bx) -- (Bxy) -- (Axy) -- cycle; % правая
  % невидимые рёбра
  \draw[dashed, thin] (A) -- (B);
  \draw[dashed, thin] (B) -- (Bx);
  \draw[dashed, thin] (B) -- (By);
  % видимые рёбра
  \draw (A) -- (Ax) -- (Axy) -- (Ay) -- cycle;
  \draw (Ay) -- (By) -- (Bxy) -- (Axy);
  \draw (Bx) -- (Bxy);
  \draw (Ax) -- (Bx);

  % ================= размеры: подписи рядом со своими рёбрами =================
  \node at (5.97,4.76)              {$\mathrm{d}x$};  % над задним верхним ребром
  \node[anchor=west] at (6.32,2.62) {$\mathrm{d}y$};  % справа от переднего правого ребра
  \node at (7.15,1.70)              {$\mathrm{d}z$};  % под нижним "глубинным" ребром

  % ================= силы давления на гранях x и x+dx =================
  % слева давление больше: стрелка длиннее
  \draw[presf] (2.03,3.35) -- (3.83,3.35)
        node[at start, left=2pt] {$P(x)\,\mathrm{d}y\,\mathrm{d}z$};
  \draw[presf] (7.93,3.35) -- (6.83,3.35)
        node[at start, right=2pt] {$P(x+\mathrm{d}x)\,\mathrm{d}y\,\mathrm{d}z$};

  % ================= потенциальная сила и её проекция на ось x =================
  \fill (5.33,3.00) circle (0.055);                    % центр бруска
  \draw[helper] (4.30,1.10) -- (4.30,3.00);            % построение проекции
  \draw[comp] (5.33,3.00) -- (4.30,3.00) node[left=2pt] {$F_x$};
  \draw[potf] (5.33,3.00) -- (4.30,1.10) node[below=3pt] {$\vec F$};

\end{tikzpicture}
\end{document}
